\begin{vidu}
	Cho c là tốc độ ánh sáng trong chân không. Theo thuyết tương đối, một hạt có khối lượng nghỉ $m_0$, khi chuyển động với tốc độ 0,8c thì có khối lượng động (khối lượng tương đối tính) là m. Biết năng lượng nghỉ của hạt là $E_0$. Tính:
	\begin{itemize}
		\item[a)] tỉ số $\dfrac {m_0}{m}$. 
		\item[b)] năng lượng toàn phần của hạt.
		\item[c)] động năng của hạt.
	\end{itemize}
\loigiai{
		\begin{itemize}
		\item[a)]  Áp dụng công thức: $m=\dfrac{m_0}{\sqrt{1-\dfrac{v^2}{c^2}}}$ với $v = 0,8c$ nên có $\dfrac {m_0}{m} =\dfrac{3}{5}$.
		\item[b)] Năng lượng toàn phần của hạt: $E=mc^2=\dfrac{1}{\sqrt{1-\dfrac{v^2}{c^2}}}m_0c^2=\dfrac{1}{\sqrt{1-0,8^2}}.E_0=\dfrac{5}{3}E_0$.
		\item[c)] Động năng của hạt: $W_{\text{đ}}=E-E_0=\left(\dfrac{5}{3}-1\right)E_0=\dfrac{2E_0}{3}$.
	\end{itemize}
}
\end{vidu}

\loai{Khối lượng tương đối tính}
\begin{vidu}%[3P7-3.2B][Loai1]
	\nguon{MH - 2017} Cho c là tốc độ ánh sáng trong chân không. Theo thuyết tương đối, một hạt có khối lượng nghỉ $m_0$, khi chuyển động với tốc độ 0,6c thì có khối lượng động (khối lượng tương đối tính) là m. Tỉ số $\dfrac {m_0}{m}$ là
	\choice
	{0,3}
	{0,6}
	{0,4}
	{\True 0,8}
	\loigiai{
		Áp dụng công thức: $m=\dfrac{m_0}{\sqrt{1-\dfrac{v^2}{c^2}}}$ với $v = 0,6 c$ nên có $\dfrac {m_0}{m} = 0,8$.
	}
\end{vidu}
\loai{Năng lượng toàn phần}
\begin{vidu}%[3P7-3.2B][Loai2]
	Kí hiệu c là vận tốc ánh sáng trong chân không. Một hạt vi mô, có năng lượng nghỉ $E_0$ và có vận tốc bằng $\dfrac{12c}{13}$ thì theo thuyết tương đối hẹp, năng lượng toàn phần của nó bằng
	\choice
	{$\dfrac{13{E_0}}{12}$}
	{$2,4{E_0}$}
	{\True $2,6{E_0}$}
	{$\dfrac{25E_0}{13}$}
	\loigiai{
		Ta có: $E=\dfrac{E_0}{\sqrt{1-\dfrac{{{(12c)}^2}}{{{13}^2}c^2}}}=\dfrac{E_0}{\sqrt{\dfrac{25}{169}}}=\dfrac{13}{5}{E_0}=2,6{E_0}$.
	}
\end{vidu}
\loai{Động năng tương đối tính}
\begin{vidu}%[3P7-3.2B][Loai3]
	\nguon{ĐH - 2010} Một hạt có khối lượng nghỉ $m_0$. Theo thuyết tương đối, động năng của hạt này khi chuyển động với tốc độ 0,6c (c là tốc độ ánh sáng trong chân không) là
	\choice
	{$1,25m_0c^2$}
	{$0,36m_0c^2$}
	{\True $0,25m_0c^2$}
	{$0,225m_0c^2$}
	\loigiai{
		\begin{itemize}
			\item Theo công thức tính động năng tương đối tính
			$K=\dfrac{m_0c^2}{\sqrt{1-\dfrac{v^2}{c^2}}}-m_0c^2$. 
			\item Thay $v = 0,6c$ vào ta được $K = 0,25m_0c^2$.
		\end{itemize}
	}
\end{vidu}
\loai{So sánh năng lượng nghỉ và động năng}
\begin{vidu}%[3P7-3.1B][Loai4]
	\nguon{CĐ - 2011} Một hạt đang chuyển động với tốc độ bằng 0,8 lần tốc độ ánh sáng trong chân không. Theo thuyết tương đối hẹp, động năng $K$ của hạt và năng lượng nghỉ  của nó liên hệ với nhau bởi hệ thức
	\choice
	{$K = \dfrac {8 E_0} {15}$}
	{$K = \dfrac {15 E_0} {8}$}
	{$K = \dfrac {3 E_0} {2}$}
	{\True $K = \dfrac {2 E_0} {3}$}
	\loigiai{
		Ta có: $K = E - E_0 = \dfrac {E_0} {\sqrt {1 - \dfrac {v^ {2}} {c^ {2}}}} - E_0 = \dfrac {E_0} {\sqrt {1 - \dfrac {0,8^ {2} \cdot c^ {2}} {c^ {2}}}} - E_0 = \dfrac {E_0} {0,6} - E_0 \Rightarrow K = \dfrac {2 E_0} {3}$.
	}
\end{vidu}
\loai{Tìm tốc độ}
\begin{vidu}%[3P7-3.2K][Loai5]
	\nguon{ĐH - 2011} Theo thuyết tương đối, một electron có động năng bằng một nửa năng lượng nghỉ của nó thì electron này chuyển động với tốc độ bằng
	\choice
	{$2,75.10^8$ m/s}
	{\True $2,24.10^8$ m/s}
	{$1,67.10^8$ m/s}
	{$2,41.10^8$ m/s}
	\loigiai{
		Ta có: $K = E - E_0 = \dfrac{1}{2}E_0\Rightarrow E=\dfrac{3}{2}E_0\Leftrightarrow \dfrac{m_0c^2}{\sqrt{1-\dfrac{v^2}{c^2}}}=\dfrac{3}{2}m_0c^2\Leftrightarrow v=2,24.10^8\ m/s$.
	}
\end{vidu}